% !TeX root = ../../Book2.tex
\chapter{Fourier 限制问题}
\label{b2:ch:23}
\BookChapterNumber{b2:ch:23}
\begin{chapterabstract}
Fourier 变换的环境空间等价类不能直接在零测度曲面上取值，限制估计为这种取值建立连续意义。本章先辨认帽集中、振荡尾部和曲率退化造成的必要条件，再借助延拓对偶、正二次型与解析族证明球面及非零曲率紧片的 Stein--Tomas 定理。选读部分继续讨论极值函数、紧性损失和现代几何方法，区分已证范围、条件性结论与仍需深层输入的目标。
\end{chapterabstract}

\begin{learninggoals}\begin{itemize}
\item\label{b2:item:23-goal-1}从稠密测试类定义曲面限制，分清普通 Lebesgue 范数控制、带矩迹估计以及不同曲面测度的作用。
\item\label{b2:item:23-goal-2}分别核对帽面积、相干管体积和远场振荡尾部，在曲率退化或测度改变时重新判断必要条件。
\item\label{b2:item:23-goal-3}把限制估计、延拓估计与卷积二次型联系起来，准确保留插值区域、半径因子及平方范数所在的空间。
\item\label{b2:item:23-goal-4}理解解析族的分布延拓、两条边界估计和目标参数，并辨认经典端点之后的紧性问题与几何问题。
\end{itemize}\end{learninggoals}

Hausdorff--Young 控制的是整个频率空间中的大小，不能直接约束一个零测度曲面上的取值。即使从 Schwartz 函数开始，这种取值也可能因频率集中在越来越薄的邻域而失控。本章首先寻找能够排除这种失控的输入空间；曲面上的限制随后由连续延拓定义，而不依赖某个任意选择的 Fourier 代表。

\ProseParagraphBreak

曲率带来两种看似相反的现象。分布在整个曲面上的光滑数据产生振荡抵消，给出远场衰减；集中在很小一顶帽内的数据，却可在一根很大的细长管中保持同相。\secref{b2:sec:2302}把这些现象转成必要条件。它们检验一个猜测是否可能成立，却不替代对不同帽相互作用的估计。

\ProseParagraphBreak

主线按限制、延拓、解析族的次序推进。对偶性把曲面取值变为一个绝对收敛的振荡积分，\(TT^*\) 再把平方限制界变为具有曲面 Fourier 支集的卷积估计。Stein--Tomas 证明把 Fourier 侧的有界性与空间侧的核控制放在不同参数边界上，插值到球面测度所在的位置。这里既要计算指数，也要解释分布族的定义及其增长，不能把一个形式参数代入发散积分。

\ProseParagraphBreak

所需基础包括 Fourier 反演与 Plancherel、Hölder 不等式、算子对偶、\chapref{b2:ch:17}的插值和\chapref{b2:ch:21}的驻相方法。选读的\secref{b2:sec:2305}与\secref{b2:sec:2306}分别讨论最佳常数之后的紧性，以及经典范围之外的几何组织；略过它们不影响本章端点定理的证明。章末前六题沿主线训练，后两题通过可显式计算的极值与 Gaussian 波包补充选读内容。

% !TeX root = ../../Book2.tex
\section{限制算子}
\label{b2:sec:2301}

Fourier 变换在 \(L^2\) 中是一个等价类，而光滑曲面在环境空间中具有零 Lebesgue 测度。因此“把 \(L^2\) Fourier 变换代入曲面”没有天然意义。限制问题问的是：从 Schwartz 函数开始的曲面取值操作，是否能由某个环境 \(L^p\) 范数控制，从而连续延拓。

\subsection{\texorpdfstring{\(\widehat f|_S\)}{Fourier 变换在 S 上的限制} 的意义}
\label{b2:subsec:230101}

对紧致光滑超曲面 \(S\subset\mathbb R^d\)、\(d\ge2\)，先定义
\[
 Rf=\widehat f|_S,\qquad f\in\mathcal S(\mathbb R^d).
\]
\symidx{\(Rf\)：曲面 Fourier 限制}
一个限制估计是
\(\|Rf\|_{L^r(S,\mathrm d\sigma)}\le C\|f\|_{L^p(\mathbb R^d)}\)。
若 \(p<\infty\)，稠密性给出唯一的 \(L^p\to L^r(S)\) 延拓。这个延拓是 \(L^r(S)\) 中的极限，并非对每个输入都先选择一个 Fourier 代表再逐点代入。

\ProseParagraphBreak
不能从 Plancherel 直接得到 \(p=r=2\)。取一个曲面坐标邻域，令
\(\widehat f_\epsilon(\xi)=b(\xi)\chi(\rho(\xi)/\epsilon)\)，其中 \(\rho=0\) 定义曲面，\(b|_S\ne0\)、\(\chi(0)=1\)。
曲面限制范数不随 \(\epsilon\) 变化，而
\(\|f_\epsilon\|_2=\|\widehat f_\epsilon\|_2=O(\epsilon^{1/2})\)。
这排除了 \(L^2\to L^2(S)\) 界，也显示环境等价类在零测集取值的困难具有真实的范数后果。

若额外要求空间矩可积，普通平方迹也可以成立。这个模型有助于把限制定理与传统迹定理区别开来。
\begin{proposition}\label{b2:prop:23-moment-trace}
对 Schwartz 输入 \(f\)，有
\[
 \|\widehat f|_{S^{d-1}}\|_2^2
 \leq C_d\|f\|_2\bigl(\|f\|_2+\||x|f\|_2\bigr).
\]
右侧有限的加权 \(L^2\) 空间因而有连续的球面迹延拓。
\end{proposition}
\begin{proof}
对每个方向 \(\omega\)，令 \(F(r)=\widehat f(r\omega)\)，在 \(1/2<r<3/2\) 上使用微积分基本定理。把起点 \(s\) 在这个区间平均，得到
\[
 |F(1)|^2\leq C\int_{1/2}^{3/2}|F(r)|^2\,\mathrm dr
 +C\int_{1/2}^{3/2}|F(r)|\,|F'(r)|\,\mathrm dr.
\]
对球面方向积分，并注意此半径区间的 \(r^{d-1}\) 与一可比。Cauchy--Schwarz 将第二项控制为
\(C\|\widehat f\|_{L^2(\text{壳})}\|\nabla\widehat f\|_{L^2(\text{壳})}\)。
Plancherel 及 \(\partial_{\xi_j}\widehat f=-2\pi\mathrm i\widehat{x_jf}\) 就给出结论。光滑紧支集逼近在带矩的平方范数中成立，估计保证其迹为 Cauchy，完成延拓。
\end{proof}
这条迹界仍是平方理论，但输入空间额外记录位置衰减或矩。Stein--Tomas 改用更小的普通 Lebesgue 指数，从而不要求输入有空间矩；它用曲面的振荡几何换取迹的控制。把上述带矩界写成无条件 \(L^2\) 限制，就会重新遇到薄壳反例，不能省去右侧的矩范数。

\subsection{曲面测度}
\label{b2:subsec:230102}

若 \(S\) 局部写成 \(\xi_d=h(\xi')\)，则
\[
 \mathrm d\sigma=\sqrt{1+|\nabla h|^2}\,\mathrm d\xi'.
\]
对光滑正密度 \(a\)，也可使用 \(\mathrm d\mu=a\,\mathrm d\sigma\)；在 \(a\) 有正上下界的紧片上，两种范数等价。若密度在某点消失，某些必要条件可能改变，所以陈述锐利估计时应明确测度。

\ProseParagraphBreak
余面积公式与频率薄壳联系紧密。对于球面，极坐标给出
\[
 \int_{\mathbb R^d}F(\xi)\,\mathrm d\xi
   =\int_0^\infty r^{d-1}\int_{S^{d-1}}F(r\omega)\,\mathrm d\sigma(\omega)\,\mathrm dr.
\]
在半径 \(1\) 的薄壳中，\(r^{d-1}\) 是有统一上下界的因子。曲面取值可以视为频率薄壳平均在厚度趋零时的极限，但没有范数控制时，这个极限仍可能发散。

\subsection{曲率}
\label{b2:subsec:230103}

本章首先以单位球面为模型，其 Gaussian 曲率处处非零。图形曲面的曲率条件可在驻点处通过 \(D^2h\) 的行列式表达；\chapref{b2:ch:21}已经解释它怎样产生 \((d-1)/2\) 个衰减幂次。
曲率对限制估计的作用并非单纯保证 Fourier 衰减。小曲面帽的法向方向还形成细长的相干管束；所有帽的叠加会产生比单个核估计复杂得多的几何问题。

\ProseParagraphBreak
平面完全不同：对 \(\xi_d=0\)，限制是输入沿 \(x_d\) 积分后在其余变量上 Fourier 变换。把输入写为 \(g(x')h(x_d/L)\)，曲面限制含一个 \(L\) 因子，而 \(L^p\) 范数只含 \(L^{1/p}\)。让 \(L\to\infty\) 排除任何 \(p>1\) 的统一界（取固定非零 \(g,h\)）。这不是一个曲率不足的小损失，而是沿平面法向可以无限相干。

\subsection{缩放}
\label{b2:subsec:230104}

令 \(S_R=RS\)，则面测度缩放为 \(R^{d-1}\mathrm d\sigma\)。
若单位曲面的 \(L^p\to L^r\) 估计成立，换元得到
\[
 \|\widehat f\|_{L^r(S_R)}\le
 C R^{(d-1)/r-d/p'}\|f\|_p.
\]
证明取 \(g(x)=f(x/R)\)，则 \(\widehat g(\xi)=R^d\widehat f(R\xi)\)，
\(\|g\|_p=R^{d/p}\|f\|_p\)，再代入面测度因子。
这一式子把半径依赖写清楚；单位球面本身不是尺度不变对象，不能把不同半径的常数都无条件设成同一个数。

% !TeX root = ../../Book2.tex
\section{必要条件}
\label{b2:sec:2302}

限制估计不是单凭 Fourier 变换的有界性就能得到的。曲面上的质量可以集中在小帽，也可以遍布整个球面；前一种输入产生很长的相干管，后一种输入留下缓慢衰减的尾部。本节分别计算这些障碍。计算所得是不等式成立的必要条件，后面的充分估计还需要独立证明。

\subsection{缩放}
\label{b2:subsec:230201}

先直接研究曲面上的积分
\[
 Eg(x)\coloneqq\int_S e^{2\pi\mathrm i x\cdot\xi}
                     g(\xi)\,d\sigma(\xi),
 \quad \|Eg\|_{L^q(\mathbb R^d)}
                 \leq C\|g\|_{L^s(S)}.
\]
\symidx{\(Eg\)：曲面延拓算子}
其中 \(S\) 暂取单位球面，\(1\leq s,q\leq\infty\)。球面测度有限，故
\(g\in L^s(S)\) 总属于 \(L^1(S)\)，积分处处绝对收敛。
这里先把 \(E\) 当作一个明确的积分算子；它与限制估计的对偶等价将在
\cref{b2:prop:23-duality}中证明。对应的限制指数是
\(p=q'\)、\(r=s'\)。

\ProseParagraphBreak
设半径 \(R\) 的球面取原生面积测度，\(g_R\) 是其上的函数，并令
\(\widetilde g(\omega)=g_R(R\omega)\)。直接换元得
\[
 E_Rg_R(x)=R^{d-1}E\widetilde g(Rx),\quad
 \|g_R\|_{L^s(RS)}
       =R^{(d-1)/s}\|\widetilde g\|_{L^s(S)}.
\]
若单位球面的最佳常数为 \(C\)，半径 \(R\) 的最佳常数便为
\[
 C_R=C R^{\,d-1-d/q-(d-1)/s}.
\]
因为输入缩放可逆，这里也是准确的最佳常数关系。
它与\cref{b2:prop:23-surface-scaling}在共轭指数下完全一致。
固定单位球面并不要求这个幂等于零；只有要求一族半径上的常数统一时，才能由 \(R\to0\) 和 \(R\to\infty\) 进一步限制该幂。
本节真正限制固定球面指数的，是球面内部任意小帽的尺度。

\ProseParagraphBreak
还须区分两个空间的测度。曲面有限测度给出
\[
 \|g\|_{s_1}\leq \sigma(S)^{1/s_1-1/s_2}\|g\|_{s_2},
 \quad 1\leq s_1\leq s_2\leq\infty.
\]
环境空间的测度却是无穷大，不能据 \(q\) 的大小直接写出一般
\(L^q(\mathbb R^d)\) 之间的包含。不过延拓函数另有
\(\|Eg\|_\infty\leq\|g\|_1\)。如果已经证明某个输出指数 \(q_0<\infty\)
的界，那么
\[
 \|Eg\|_q\leq
 \|Eg\|_{q_0}^{q_0/q}
       \cdot\|Eg\|_\infty^{\,1-q_0/q},
 \quad q_0\leq q<\infty,
\]
才给出更大输出指数的界。这一步使用的是输出的额外有界性。

\subsection{Knapp 例子}
\label{b2:subsec:230202}

球面小帽在切向上的半径为 \(\delta\)，离开切平面的高度只有
\(O(\delta^2)\)。为了使不同频率的相位保持接近，环境变量的切向分量可以达到
\(\delta^{-1}\)，法向分量则可以达到 \(\delta^{-2}\)。
这就是下面的\term[knapp-lizi]{Knapp 例子}[Knapp example]所用的几何。

\Needspace{6\baselineskip}
\begin{proposition}[帽与管的必要条件]\label{b2:prop:23-knapp}
设 \(d\geq2\)，单位球面 \(S^{d-1}\) 取面积测度，
\(1\leq s,q\leq\infty\)。若其延拓积分对全部
\(g\in L^s(S^{d-1})\) 满足
\(\|Eg\|_q\leq C\|g\|_s\)，则必须有
\[
 \frac{d+1}{q}\leq\frac{d-1}{s'},
 \quad \frac1s+\frac1{s'}=1.
\]
\end{proposition}
\begin{proof}
在北极附近写
\(\xi=(\eta,h(\eta))\)，其中
\(h(\eta)=\sqrt{1-|\eta|^2}\)。取 \(0<\delta<1/4\)，令
\[
 \Theta_\delta=\{(\eta,h(\eta)):|\eta|\leq\delta\},
 \quad g_\delta=\mathbf1_{\Theta_\delta}.
\]
面积因子 \((1-|\eta|^2)^{-1/2}\) 在这些帽上有统一正上下界，所以
\[
 \sigma(\Theta_\delta)\asymp\delta^{d-1},\quad
 \|g_\delta\|_s\asymp\delta^{(d-1)/s}.
\]
当 \(s=\infty\) 时后一范数等于一，仍与所写幂一致。

由 \(h(\eta)-1=-|\eta|^2/2+O(|\eta|^4)\)，存在与 \(\delta\)
无关的 \(C_0\)，使 \(|h(\eta)-1|\leq C_0\delta^2\)。
取一个固定的小 \(c>0\)，在集合
\[
 \mathcal T_\delta=
 \{x:|x'|\leq c\delta^{-1},\ |x_d|\leq c\delta^{-2}\}
\]
上，有
\[
 2\pi|x'\cdot\eta+x_d(h(\eta)-1)|
       \leq 2\pi c(1+C_0).
\]
选 \(c\) 使右侧不超过 \(\pi/3\)。从积分中提出公共因子
\(e^{2\pi\mathrm i x_d}\)，余下被积函数的实部至少为 \(1/2\)，故
\[
 |Eg_\delta(x)|\geq\tfrac12\sigma(\Theta_\delta)
       \gtrsim\delta^{d-1}\quad(x\in\mathcal T_\delta).
\]
管的体积为常数乘 \(\delta^{-(d+1)}\)。因此，包括 \(q=\infty\)
的通常解释在内，
\[
 \|Eg_\delta\|_q\gtrsim
       \delta^{d-1-(d+1)/q}.
\]
假定的估计会要求
\(\delta^{d-1-(d+1)/q}\lesssim\delta^{(d-1)/s}\)。
令 \(\delta\downarrow0\)，左侧幂不能小于右侧幂，整理即得结论。
\end{proof}

同样的构造可以用非负光滑帽代替特征函数：令帽内部保留一个固定比例的核心，振幅在核心上有正下界即可。因此，必要条件并非由输入的不光滑造成。对任意光滑曲面的局部片，减去切平面后只需使用高度的 \(O(\delta^2)\) 上界，也可得到同一个必要条件；若曲率退化，还可能有更强的条件。

\ProseParagraphBreak
当 \(s=2\) 时，上式成为
\[
 q\geq q_*\coloneqq\frac{2(d+1)}{d-1}.
\]
本章将证明等号成立，因此平方曲面输入的最小输出指数已经由一个帽确定。
例如 \(d=3\) 时帽面积约为 \(\delta^2\)，管体积约为 \(\delta^{-4}\)，
比较 \(\delta^{2-4/q}\) 与 \(\delta\) 就得到 \(q\geq4\)。
若 \(s=1\)，必要条件要求 \(q=\infty\)，正好对应绝对积分界。

\cref{b2:fig:23-cap-tube}把两个小相位条件画成空间宽度与长度。两幅图分别属于频率空间和位置空间；帽的法向高度与管的长度不是同一空间中的两个长度。

\begin{figure}[htbp]
\centering
\begin{tikzpicture}[x=1cm,y=1cm]
\begin{scope}[shift={(-2.8,0)}]
\draw[->] (-1.5,0)--(1.5,0) node[right]{\(\eta\)};
\draw[->] (0,-.3)--(0,1.8) node[above]{法向};
\draw[AnalysisBlue,thick,domain=-1.3:1.3,samples=35] plot(\x,{1.3-.35*\x*\x});
\draw[AnalysisBlue,very thick,domain=-.5:.5,samples=15] plot(\x,{1.3-.35*\x*\x});
\draw[<->] (-.5,.7)--(.5,.7) node[midway,below]{\(\delta\)};
\draw[<->] (.7,1.2125)--(.7,1.3) node[right]{\(\delta^2\)};
\node at (0,-.6){频率帽：切向与法向高度};
\end{scope}
\draw[->] (-.7,.8)--(.6,.8);
\begin{scope}[shift={(2.8,0)}]
\fill[AnalysisBlue!15] (-.4,0) rectangle(.4,1.8);
\draw[AnalysisBlue,thick] (-.4,0) rectangle(.4,1.8);
\draw[<->] (-.4,-.2)--(.4,-.2) node[midway,below]{\(\delta^{-1}\)};
\draw[<->] (.65,0)--(.65,1.8) node[midway,right]{\(\delta^{-2}\)};
\node at (0,2.05){\(x_d\) 法向};
\node at (0,-.85){物理空间相干管};
\end{scope}
\end{tikzpicture}
\caption[频率帽与相干管]{Knapp 尺度的对应，不按比例。移去北极的常相位后，要让 \(x'\cdot\eta\) 与 \(x_dO(|\eta|^2)\) 都小，分别需要 \(|x'|\lesssim\delta^{-1}\) 与 \(|x_d|\lesssim\delta^{-2}\)。细长管的长轴沿帽中心的法向；帽的面积和管的体积仍由正文计算。}
\label{b2:fig:23-cap-tube}
\end{figure}

\subsection{球面尾部的可积性条件}
\label{b2:subsec:230203}

小帽检验只利用一个大而有限的管。要检查无穷远处的积分，取一个不会随尺度收缩的输入更合适。球面上的常数函数便提供了这样的检验。

\Needspace{6\baselineskip}
\begin{proposition}[常密度的严格尾部条件]\label{b2:prop:23-sphere-tail}
设 \(d\geq2\)，\(E\) 是单位球面面积测度的延拓算子。
对于 \(1\leq q<\infty\)，有
\[
 E1\in L^q(\mathbb R^d)
 \quad\Longleftrightarrow\quad
 q>\frac{2d}{d-1}.
\]
因而任何 \(L^s(S^{d-1})\to L^q(\mathbb R^d)\) 的有界性，
若 \(q<\infty\)，都必须满足这个严格不等式。
\end{proposition}
\begin{proof}
球面的旋转不变性使 \(E1(r\omega)\) 只依赖 \(r\)。
相位 \(\omega\cdot\xi\) 在球面上的驻点恰为 \(\omega\) 与 \(-\omega\)，
两处的切向 Hessian 都非退化。按\cref{b2:thm:21-stationary-phase}
分别取两个小驻点邻域，其余部分用非驻相分部积分，可得
\[
 E1(r\omega)=r^{-(d-1)/2}
       \bigl(c_+e^{2\pi\mathrm i r}+c_-e^{-2\pi\mathrm i r}\bigr)
       +O(r^{-(d+1)/2}),\quad r\geq1.
\]
系数 \(c_\pm\) 非零；由球面对称性，它们互为共轭，因而模相等。
旋转不变性也使余项常数不依赖 \(\omega\)。
这给出 \(|E1(r\omega)|\leq Cr^{-(d-1)/2}\)。
原积分在有界区域有界，所以当
\(d-1-q(d-1)/2<-1\) 时，极坐标积分收敛。

反过来，令
\(H(r)=c_+e^{2\pi\mathrm i r}+c_-e^{-2\pi\mathrm i r}\)。
这是非恒零的一周期连续函数。因此存在 \(\kappa>0\) 和
\(A\subset[0,1]\)，使 \(|A|>0\)，且 \(|H(t)|\geq2\kappa\)
对 \(t\in A\) 成立。对充分大的整数 \(n\)，余项在
\(n+A\) 上至多为 \(\kappa r^{-(d-1)/2}\)，从而
\[
 |E1(r\omega)|\geq\kappa r^{-(d-1)/2}
       \quad(r\in n+A).
\]
于是
\[
 \|E1\|_q^q\gtrsim
 \sum_{n\geq n_0} n^{\,d-1-q(d-1)/2}.
\]
这个级数只有在指数小于 \(-1\) 时才收敛。等号时出现调和级数，
也发散。常数函数属于球面上的每个 \(L^s\)，故得到最后的必要条件。
\end{proof}

这里不能把“驻相衰减上界”直接当作发散证明；真正排除端点的是每段径向区间内都有固定比例的下界集合。余弦式领先项确有零点，但这些零点无法消去其余径向区域的质量。

\ProseParagraphBreak
两种必要条件的作用不同：常密度给出严格的水平门槛
\(q>2d/(d-1)\)，小帽给出可含等号的关系
\((d+1)/q\leq(d-1)/s'\)。
在 \(s=2\) 时，小帽条件更强；若允许更大的 \(s\)，两个边界便可能分别起作用。
三维中它们是 \(q>3\) 和 \(4/q\leq2/s'\)。

\subsection{几何障碍}
\label{b2:subsec:230204}

曲率退化时，圆形帽未必反映最强的集中机制。考虑三维中的局部曲面
\[
 S=\{(u,v,u^2+v^4):|u|,|v|<a\},
\]
其中 \(a>0\) 固定，取原生面积测度。对小 \(\delta>0\)，取帽
\[
 |u|\leq\delta^{1/2},\quad |v|\leq\delta^{1/4}.
\]
其高度至多为 \(2\delta\)，面积约为 \(\delta^{3/4}\)。
令 \(g_\delta\) 为该帽的特征函数。若
\[
 |x_1|\leq c\delta^{-1/2},\quad
 |x_2|\leq c\delta^{-1/4},\quad
 |x_3|\leq c\delta^{-1},
\]
则相位 \(2\pi(x_1u+x_2v+x_3(u^2+v^4))\) 的绝对值至多为
\(8\pi c\)。取 \(c\) 足够小，仍可用实部正下界得到
\(|Eg_\delta|\gtrsim\delta^{3/4}\)。
这个长方管的体积约为 \(\delta^{-7/4}\)，所以有界性必须要求
\[
 \delta^{3/4-7/(4q)}
       \lesssim\delta^{3/(4s)},\quad
 \frac7q\leq\frac3{s'}.
\]
特别地，平方输入必须有 \(q\geq14/3\)，比非零曲率曲面的三维指数
\(4\) 更严格。这里曲率在 \(v=0\) 处消失，第四次方方向让帽比二次帽更宽，才产生了更大的相干体积。

\ProseParagraphBreak
这些检验也说明了必要条件的限度。它们都选取了一类特别的输入；通过这些检验，只能说明该类输入没有排除某个指数。任意输入还可能包含许多尺度、位置与方向不同的帽，叠加后的管束相互作用没有包含在单帽计算中。下一节建立对偶与 \(TT^*\) 表述，随后用解析插值证明平方输入的临界估计；更广的目标放在选读部分讨论。

% !TeX root = ../../Book2.tex
\section{延拓算子}
\label{b2:sec:2303}

延拓算子把曲面数据变成环境空间的振荡函数。对偶性把无法直接解释的曲面取值，转化为一个处处有意义的积分，并为 \(TT^*\) 方法提供起点。

\subsection{对偶限制}
\label{b2:subsec:230301}

\begin{proposition}\label{b2:prop:23-duality}
对于 \(1\le p,r<\infty\)，估计 \(R:L^p\to L^r(S)\) 等价于
\(E:L^{r'}(S)\to L^{p'}(\mathbb R^d)\)，两者最佳常数相同。
\end{proposition}
\begin{proof}
对 Schwartz 函数与曲面光滑函数，Fubini 给出
\[
 \int_S\widehat f(\xi)\overline{g(\xi)}\,\mathrm d\sigma
 =\int_{\mathbb R^d} f(x)\overline{Eg(x)}\,\mathrm dx.
\]
一侧估计配以 Hölder 不等式，再对另一侧单位球取上确界，得到反向算子的估计。有限指数时稠密性完成延拓；出现无穷对偶指数时使用通常的 \(L^1\) 对偶 \(L^\infty\) 即可。最佳常数由两次同样操作相等。
\end{proof}
由于曲面有限测度，\(Eg\) 对 \(g\in L^1(S)\) 绝对收敛且有界连续。若 \(g\) 更粗，先在已建立的范数估计下通过稠密极限定义输出；不能把所有延拓数据都假定光滑。

\subsection{Fourier 衰减}
\label{b2:subsec:230302}

光滑曲面密度的 Fourier 衰减由\chapref{b2:ch:21}的驻相公式给出。在非零曲率紧片上，
\(|\widehat{\mathrm d\mu}(x)|\lesssim(1+|x|)^{-(d-1)/2}\)。
若 \(g\) 光滑，则同样结论成立，但常数含有 \(g\) 的导数。\(L^2(S)\) 单位球不控制这些导数，因此不能把这个逐输入衰减直接当作 \(L^2\to L^q\) 算子界。

\ProseParagraphBreak
平移环境变量对应曲面调制：
\[
 E(e^{-2\pi\mathrm ia\cdot\xi}g)(x)=Eg(x-a).
\]
于是固定 \(L^2\) 范数的输入可使输出峰值移到任意远处，排除一个仅由 \(\|g\|_2\) 控制、固定以原点为中心的逐点衰减界。算子估计必须容许这种平移对称性。

\subsection{\texorpdfstring{\(L^p\to L^q\)}{Lp→Lq} 估计}
\label{b2:subsec:230303}

显然 \(\|Eg\|_\infty\le\|g\|_1\)。若已经得到
\(\|Eg\|_{q_0}\le C\|g\|_2\)，曲面有限测度还给出
\(\|Eg\|_\infty\le |S|^{1/2}\|g\|_2\)。对输出两个范数插值，可得到每个 \(q\ge q_0\) 上的 \(L^2\to L^q\) 界。
对输入指数也可与 \(L^1\to L^\infty\) 插值，产生一个三角范围；这些是已证估计的推论，不能据此覆盖全部必要条件范围。

\ProseParagraphBreak
在球面上，曲面函数 \(g=1\) 的尾部检验已经说明输出不属于任意小的 \(q\)。紧曲面只在输入端是有限测度空间，输出端仍然在整个 \(\mathbb R^d\)。将两端都按有限测度嵌入处理会得到错误的范围。

% 扩写定位：b2:subsec:230303
一旦建立 \(L^2(S)\to L^{q_*}\) 的延拓估计，基本插值就给出一片完整的指数区域。另一方面，曲面测度有限，
\(\|Eg\|_\infty\leq\|g\|_1\)。对 \(0\leq\theta\leq1\)，在这两个端点间使用 Riesz--Thorin，得到
\[
 \frac1s=1-\frac\theta2,\qquad
 \frac1q=\frac\theta{q_*},\qquad
 \|Eg\|_q\leq C_\theta\|g\|_s.
\]
故对 \(1\leq s\leq2\)，证明了直线
\(1/q=\frac{d-1}{d+1}(1-1/s)\) 上的估计。再与同一输入指数的 \(L^\infty\) 界插值，可增大输出 \(q\)。对 \(s>2\)，有限曲面测度使 \(L^s(S)\subset L^2(S)\)，于是仍得到所有 \(q\geq q_*\) 的估计。这后一范围来自输入空间包含，不是把刚才的直线延长到 \(\theta>1\)。

\ProseParagraphBreak

对偶地，从 \(L^{p_*}(\mathbb R^d)\to L^2(S)\) 与 \(L^1(\mathbb R^d)\to L^\infty(S)\) 得到
\[
 \frac1p=1-\frac{\theta}{q_*},\qquad \frac1r=\frac\theta2.
\]
这使一个端点定理同时控制较小输入指数和较大曲面输出指数。曲面测度有限，可以把输出指数降低；环境空间测度无穷，却不能据此任意提高输入 \(p\)。这两个方向在实际引用时很容易混淆，必须先区分哪一侧是有限测度空间。

\ProseParagraphBreak

上述区域仍然只是一个已证明的充分区域。Knapp 管的必要条件对更一般 \(s,q\) 仍适用，径向尾部则给出另外的严格限制；两者之间是否还有空缺，需要新的估计。把插值区域画到必要条件所在边界并不会填补该空缺。指数图的作用是帮助核对已知界，而不是替代其证明。

\subsection{\texorpdfstring{\(TT^*\)}{TT*} 表述}
\label{b2:subsec:230304}

在测试函数上，
\[
 ERf=f*\check\sigma,\qquad
 \|Rf\|_2^2=\langle f*\check\sigma,f\rangle.
\]
因此若卷积算子 \(T=f*\check\sigma\) 有界 \(L^p\to L^{p'}\)，就有
\(\|Rf\|_2\le\|T\|^{1/2}\|f\|_p\)。
反过来，限制界与其对偶延拓界复合也给出同一个卷积界。这里 \(TT^*\) 的实质是一个正的二次型，即便卷积核本身不是非负函数。

\ProseParagraphBreak
若只取核绝对值并使用 Hardy–Littlewood–Sobolev，核的衰减阶 \((d-1)/2\) 会给出较差的指数。Stein–Tomas 的改进保留核的振荡和 Fourier 侧的曲面结构。下一节用解析族把这两类信息放到不同边界线上，再插值合并。

球面延拓还有一个简单的偏微分方程解释。对 \(g\in L^2(S^{d-1})\)，有限曲面测度使 \(g\in L^1(S)\)，因此可以逐阶对 \(x\) 求导，并有全部导数的统一上界。直接计算得到
\[
 (-\Delta-4\pi^2)Eg=0.
\]
这是一族光滑的 Helmholtz 解，频率恰好集中于单位能量面。光滑性来自频率支集紧致，而衰减及全空间可积性则需要曲率和输入性质；两者不可混同。

\ProseParagraphBreak

任何非零的这种解都不属于全空间 \(L^2\)。若 \(Eg\in L^2\)，其 Fourier 变换由 Plancherel 为一个 \(L^2\) 密度；另一方面，温和分布 Fourier 计算给出 \(g\,\mathrm d\sigma\)，支于 Lebesgue 零测的球面。一个绝对连续平方密度若只支于这个零测集只能为零，故 \(g=0\)。如果 \(Eg\in L^q\)、\(q<2\)，又因它有界，插值也会给出 \(L^2\)，同样不可能。限制估计的对偶输出因而自然位于大于二的指数，而不是一个通常的平方能量空间。

\ProseParagraphBreak

这也解释了 \(TT^*\) 中的平方空间究竟在哪里：平方范数落在曲面数据上，环境输出使用较大的 \(q\)。将两者都误写成全空间 \(L^2\)，会把曲面测度与 Lebesgue 测度之间的区别抹掉。用于散射和谱问题时，球面延拓是一种固定能量的广义波，而不是一个具有有限全空间平方质量的束缚态。

% !TeX root = ../../Book2.tex
\section{Stein–Tomas 定理}
\label{b2:sec:2304}

单位球面的锐利 \(L^2\) 输入限制范围可由一个显式解析族证明。下面同时说明核的估计为何在复参数下仍然有效，避免把只针对实阶 Bessel 函数的渐近式直接用于一条复竖线。

\subsection{球面测度的 Fourier 衰减}
\label{b2:subsec:230401}

由\cref{b2:thm:A2-sphere-fourier}，采用\cref{b2:def:A2-bessel-j}的 Bessel 函数归一化，有
\[
 \check\sigma(x)=2\pi |x|^{1-d/2}J_{d/2-1}(2\pi|x|).
\]
原点处按连续极限解释。由\cref{b2:thm:A2-bessel-asymptotic}的实阶渐近或\chapref{b2:ch:21}直接驻相，得到
\(|\check\sigma(x)|\lesssim(1+|x|)^{-(d-1)/2}\)。
这个指数取决于曲面维数，而不是环境维数的一半；球面有 \(d-1\) 个切向振荡变量。

\ProseParagraphBreak
在小 \(r\) 时，\(J_\nu(r)\sim(r/2)^\nu/\Gamma(\nu+1)\)，故上式有有限极限，等于球面面积。大 \(r\) 时则有两个相反法向驻点的贡献。将小尺度与大尺度分开，是后面的解析族核界能够在整个空间成立的原因。

\subsection{\texorpdfstring{\(TT^*\)}{TT*}}
\label{b2:subsec:230402}

定义分布解析族
\[
 G_z(\xi)=\frac{(1-|\xi|^2)_+^z}{\Gamma(z+1)},\qquad
 T_zf=\mathcal F^{-1}(G_z\widehat f).
\]
先对 \(\operatorname{Re}z>-1\) 定义积分，再在径向变量 \(t=1-r^2\) 中反复分部积分延拓：\(t_+^z/\Gamma(z+1)\) 的分布导数是
\(t_+^{z-1}/\Gamma(z)\)，所以没有 Gamma 极点造成的障碍。
在 \(z=-1\) 时 \(G_{-1}=\delta(1-|\xi|^2)=\tfrac12\mathrm d\sigma\)。
对 \(\operatorname{Re}z=0\)，Plancherel 给出
\[
 \|T_{it}\|_{2\to2}\le|\Gamma(1+it)|^{-1}.
\]
\cref{b2:thm:A2-bochner-riesz}的积分与幂级数计算在复参数的收敛半平面中同样成立，给出
\[
 \check G_z(x)=\pi^{-z}|x|^{-d/2-z}J_{d/2+z}(2\pi|x|).
\]
两侧先在积分收敛的半平面相等，再由分布解析性延拓。

\ProseParagraphBreak
令 \(a=(d+1)/2\)。在 \(\operatorname{Re}z=-a\) 上，Bessel 阶为
\(\nu=-1/2+it\)。我们需要
\(\sup_x|\check G_{-a+it}(x)|\le C e^{C(1+t^2)}\)。
对 \(0<r\le1\)，使用幂级数
\[
 r^{-\nu}J_\nu(2\pi r)
 =\pi^\nu\sum_{k\ge0}\frac{(-\pi^2r^2)^k}{k!\Gamma(k+\nu+1)}.
\]
因 \(|\ell+1/2+it|\ge\ell+1/2\)，Gamma 递推控制整个级数及其 \(r\) 导数。Weierstrass 乘积中
\(\log(1+t^2/n^2)\le t^2/n^2\) 给出
\(|\Gamma(1/2+it)|^{-1}+|\Gamma(1+it)|^{-1}\le C e^{C(1+t^2)}\)。
对 \(r\ge1\)，令 \(w(r)=r^{1/2}J_\nu(2\pi r)\)，Bessel 方程化为
\[
 w''+(2\pi)^2w=(\nu^2-\tfrac14)r^{-2}w.
\]
用正能量 \(|w'|^2+(2\pi)^2|w|^2\) 求导并积分，Gronwall 不等式与
\(\int_1^\infty r^{-2}\,\mathrm dr=1\) 控制 \(w,w'\)，常数至多再乘 \(e^{C(1+t^2)}\)。
因此 \(|J_\nu(2\pi r)|\le C e^{C(1+t^2)}r^{-1/2}\)；与前面的 \(r^{1/2}\) 因子相消，得到所需核界。

\subsection{Stein–Tomas 定理}
\label{b2:subsec:230403}

下面采用\term[stein-tomas-dingli]{Stein--Tomas 定理}[Stein--Tomas theorem]的下述版本。

\begin{theorem}[Stein--Tomas 定理：球面]\label{b2:thm:23-stein-tomas}
设 \(d\ge2\)、\(S=S^{d-1}\)。若
\[
 1\le p\le p_*=\frac{2(d+1)}{d+3},
\]
则 \(\|\widehat f|_S\|_2\le C_{d,p}\|f\|_p\)。
对偶地，\(E:L^2(S)\to L^q(\mathbb R^d)\) 对全部
\(q\ge q_*=2(d+1)/(d-1)\) 有界。
\end{theorem}
\begin{proof}
在带 \(-a\le\operatorname{Re}z\le0\) 中对 \(T_z\) 使用\chapref{b2:ch:17}解析族插值。
先乘标量因子 \(e^{\kappa z^2}\)，选 \(\kappa\) 大于上述竖线增长常数，使两侧算子范数在虚部方向有界；在固定实插值点这个因子只是一个有限非零常数。
右边界有 \(L^2\to L^2\) 界，左边界由有界卷积核有 \(L^1\to L^\infty\) 界。
目标 \(z=-1\) 从右边界到左边界的比例为
\(\theta=1/a=2/(d+1)\)。插值给出
\[
 \frac1p=\frac{1-\theta}{2}+\theta
       =\frac{d+3}{2(d+1)},\qquad
 \frac1q=\frac{1-\theta}{2}=\frac{d-1}{2(d+1)}.
\]
因此 \(T_{-1}=\tfrac12 ER\) 有界 \(L^{p_*}\to L^{p_*'}\)。
上一节的正二次型给出端点限制界。与显然的
\(L^1\to L^2(S)\) 界插值得到 \(1\le p\le p_*\)。
对偶及延拓输出的 \(L^\infty\) 界给出全部 \(q\ge q_*\)。
解析配对先在 Schwartz 函数上定义；分布族的增长经同一标量因子控制，故解析族插值的适用条件已经核验。
\end{proof}
Knapp 例子说明在 \(L^2(S)\) 输入这一条线上，\(q_*\) 不能减小。这是针对固定输入指数的锐利性，不意味着整个限制问题已解决。

% 扩写定位：b2:subsec:230403
解析族的分布意义也应当在整条带上说明。以法向变量 \(u=1-|\xi|^2\) 为例，基本的一维族
\(u_+^z/\Gamma(z+1)\) 最初只在 \(\operatorname{Re}z>-1\) 定义。对测试函数 \(F\) 在零点作 \(N-1\) 阶 Taylor 减法，零点附近的配对写为
\[
 \frac1{\Gamma(z+1)}
 \left[\int_0^1u^z
 \left(F(u)-\sum_{k=0}^{N-1}\frac{F^{(k)}(0)}{k!}u^k\right)\mathrm du
 +\sum_{k=0}^{N-1}\frac{F^{(k)}(0)}{k!(z+k+1)}\right].
\]
积分在 \(\operatorname{Re}z>-N-1\) 中解析；显示的简单极点被 \(1/\Gamma(z+1)\) 在负整数的零点消去。零点以外的部分没有局部可积性障碍。增大 \(N\) 就得到整个复平面上的分布解析延拓。

\ProseParagraphBreak

在 \(z=-1\) 时，只有 \(k=0\) 的极点贡献非零极限，因而得到 \(\delta_0\)；换到球的法向坐标时，正的测度因子为 \(1/|\nabla(1-|\xi|^2)|=1/2\)，所以 \(G_{-1}=\sigma/2\)。这解释了端点证明中的二分之一，不能把它吸收后仍声称得到了原归一化的最佳常数。

\ProseParagraphBreak

Taylor 减法还给出带上的增长界。在一个固定有限实部区间中，选定足够大的 \(N\)，积分由测试函数的有限个半范数控制，分母 \(z+k+1\) 的奇点与 Gamma 零点共同解释为解析函数。前面 Gamma 乘积给出的界允许至多 \(e^{C(1+|\operatorname{Im}z|^2)}\) 增长。于是乘上 \(e^{\kappa z^2}\)，选 \(\kappa>C\)，在竖直方向得到足够的衰减。对 Schwartz 输入输出的配对，所得族满足三线引理的增长条件。这个检查补足了仅验证两条边界算子范数还不足以自动进行解析插值的一步。

\ProseParagraphBreak

对一般非零曲率紧致曲面，其 \(L^2\) 延拓端点与球面一致，但上述径向解析族不能原样使用。局部法向分布、振幅和曲率需要重新组织。正式发表的 Bourgain--Demeter 论文\cite{BourgainDemeter2015Proof}的 Theorem 2.1 也明确列出这项经典端点作为背景工具；本节一般曲面的证明提纲保持深层结论状态，而完整球面证明已经独立给出。

\subsection{非零曲率超曲面}
\label{b2:subsec:230404}

\begin{theorem}[深层结构定理：非零曲率超曲面的 Stein--Tomas 估计]\label{b2:thm:23-curved-tomas}
设 \(S\) 是光滑超曲面，\(a\,\mathrm d\sigma\) 为紧支撑光滑密度，且支撑上 Gaussian 曲率处处非零。则同样的
\(L^{p_*}(\mathbb R^d)\to L^2(S,a\,\mathrm d\sigma)\) 估计成立（正密度情形）。
\end{theorem}
\begin{proofsketch}
在小片上用图形坐标，构造沿法向的分布解析族，使目标参数成为曲面测度。一个边界由符号的有界性给出 \(L^2\) 界；另一个边界对法向变量作 Fourier 变换，再在 \(d-1\) 个切向变量中用一致驻相估计，给出有界核。插值与正二次型完成证明，有限分割再拼接各片。
本节完整证明了球面的显式族；一般曲面中法向分布族、切向驻相余项的联合参数控制不在此展开，不能仅凭实参数核衰减代替这一环节。完整一般理论见 \cite{Stein1993Harmonic,Sogge2017Fourier}。
\end{proofsketch}
这里的选取是有意的：核心球面模型不依赖这个深层输入，后面的球面 Bochner–Riesz 讨论也只使用已完整证明的球面定理。一般曲面的陈述用于界定理论范围，并不作为未证明的隐含前置。

% !TeX root = ../../Book2.tex
\OptionalSection{极值问题}{b2:sec:2305}

算子有界以后，可以进一步研究最佳常数及达到它的输入。极值问题比有界性多出一个紧性要求；对称群常使普通的紧性失效。本节给出可核验的框架，不宣称所有曲面和指数都存在极值函数。

\subsection{最佳常数}
\label{b2:subsec:230501}

固定一个已证的界 \(E:L^2(S)\to L^q\)、\(q<\infty\)，令
\[
 \mathcal C=\sup_{\|g\|_2=1}\|Eg\|_q.
\]
当测度非零时 \(\mathcal C>0\)，有界性使其有限。若存在单位向量达到上确界，称为极值函数。最佳常数的数值、是否达到及所有极值函数的分类是三个不同问题，不应混成一个陈述。

\ProseParagraphBreak
即使算子范数有限，也未必有极值向量。例如 \(\ell^2\) 上的对角算子
\((x_n)\mapsto((1-1/n)x_n)\)（从 \(n=2\) 开始）范数为 \(1\)，但每个非零向量的范数都严格变小。单位坐标向量构成极值化序列，却没有强收敛子列。这个模型提醒我们必须寻找实际的紧性机制。

\subsection{极值化序列}
\label{b2:subsec:230502}

按上确界定义，可选 \(\|g_n\|_2=1\)、\(\|Eg_n\|_q\to\mathcal C\)。
Hilbert 空间弱紧性提供弱收敛子列，但弱极限可能为零。对于固定 \(x\)，曲面指数函数属于 \(L^2(S)\)，所以 \(g_n\rightharpoonup g\) 蕴含 \(Eg_n(x)\to Eg(x)\)；然而逐点收敛仍不保证 \(L^q\) 强收敛。

\ProseParagraphBreak
一个方便的充分判据是：若经对称归一化后 \(g_n\to g\) 在 \(L^2(S)\) 中强收敛，则连续性给出
\(\|g\|_2=1\)、\(\|Eg\|_q=\mathcal C\)。
该判据的证明很短，困难全在证明归一化后的强紧性，而不是上确界的形式定义。

\subsection{对称性}
\label{b2:subsec:230503}

球面的正交旋转保持输入及输出范数；乘一个常数相位也保持范数。环境平移由曲面调制实现，所以
\(g_n(\xi)=e^{-2\pi\mathrm ia_n\cdot\xi}g(\xi)\) 在 \(|a_n|\to\infty\) 时可使输出逃向无穷远。
在曲面局部放大极限中，还可能出现小帽集中，极限几何是切平面附近的抛物面；这不是单位球面精确的缩放对称，却是紧性分析中的真实损失机制。

\ProseParagraphBreak
归一化应分别处理精确对称与逼近尺度。除以相位、旋转、调制并不自动排除帽集中；排除帽集中也不自动排除多个远离的输出峰。每一种紧性损失都需要独立估计。

\subsection{极值函数}
\label{b2:subsec:230504}

若极值函数存在，在适当可微的范围 \(q>2\) 下，对变分 \(g+\epsilon h\) 求导可得到 Euler–Lagrange 方程
\[
 E^*(|Eg|^{q-2}Eg)=\mathcal C^qg
 \qquad(\|g\|_2=1).
\]
左侧按 \(L^{q'}\to L^2(S)\) 的伴随估计解释；因为
\(|Eg|^{q-2}Eg\in L^{q'}\)，该表达确实属于 \(L^2(S)\)。
求导时对复函数分别取实部方向和虚部方向，才能得到完整的复 Hilbert 空间等式。

\ProseParagraphBreak
这个方程是必要条件，不能反过来把任意解都称为全局极值函数。它还可能有局部临界点或鞍点；确认最佳常数仍需比较整个单位球。对于特殊曲面和特殊指数，卷积几何可能使常数与极值分类可显式计算，但本章不把这类特殊结论推广到一般情形。

\subsection{集中紧性}
\label{b2:subsec:230505}

集中紧性方法先寻找有正量的局部输出，再通过对称群把它移回固定区域；随后分离一个非零轮廓与余项，并用范数解耦判断余项能否继续携带极值量。\appref{b2:app:F}将说明消失、分裂和紧性三种模式，以及它们在变分问题中的作用。

\ProseParagraphBreak
一个典型的排除分裂原理是 \(q>2\) 时
\(a^{q/2}+(1-a)^{q/2}<1\)、\(0<a<1\)。
若两个轮廓在输出端充分分离、输入平方范数分别为 \(a\) 与 \(1-a\)，则两输出的 \(q\) 次范数之和至多
\(\mathcal C^q[a^{q/2}+(1-a)^{q/2}]\)，严格小于最优值。
使用这项原理之前必须证明轮廓解耦与没有其他损失；仅写这条数值不等式不能替代集中紧性的完整证明。

% 扩写定位：b2:subsec:230505
有一个可以完整证明的条件紧性判据，它把“输出没有逃走”与“输入达到极值”直接连接起来。设 \(S\) 紧致，\(E:L^2(S)\to L^q(\mathbb R^d)\) 有界，\(2<q<\infty\)，并有极值化序列 \(\|g_n\|_2=1\)。若在适当对称归一化后，
\[
 \lim_{R\to\infty}\sup_n\int_{|x|>R}|Eg_n(x)|^q\,\mathrm dx=0,
\]
则它具有强 \(L^2(S)\) 收敛子列，其极限为极值函数。

\ProseParagraphBreak

证明如下。取弱收敛子列 \(g_n\rightharpoonup g\)。由于 \(S\) 紧致，各阶导数都满足
\[
 |\partial_x^\alpha Eg_n(x)|
 \leq(2\pi)^{|\alpha|}
 \left(\int_S|\xi^\alpha|^2\,\mathrm d\sigma(\xi)\right)^{1/2}.
\]
因而输出在每个紧球上统一有界且等度连续。弱收敛给出逐点极限 \(Eg\)，Arzelà--Ascoli 及有限覆盖遂使 \(Eg_n\to Eg\) 在每个紧球上一致收敛。假设的统一尾部小性与 Fatou 控制极限尾部，故全空间 \(L^q\) 中也强收敛。于是 \(\|Eg\|_q=\mathcal C\)。弱下半连续性给出 \(\|g\|_2\leq1\)，而算子最佳界给出 \(\mathcal C\leq\mathcal C\|g\|_2\)；在非零算子情形 \(\mathcal C>0\)，所以 \(\|g\|_2=1\)。Hilbert 空间中弱收敛加范数收敛意味着强收敛，命题得证。

\ProseParagraphBreak

这一判据的假设很强，也明确指出了还需证明什么。调制产生的平移会破坏统一尾部；帽集中可能产生尺度越来越大的管，也可能破坏尾部。仅证明每个固定 \(n\) 的 \(Eg_n\) 可积，并不能把其尾部界变成对全部 \(n\) 一致的小量。因此上面的命题没有无条件断言一般曲面的极值函数存在，而是为具体问题提供一个可核对的紧性目标。

% !TeX root = ../../Book2.tex
\OptionalSection{现代 Fourier 限制问题}{b2:sec:2306}

本节区分已完整证明的 Stein–Tomas 范围、必要条件给出的广泛目标和现代几何方法。研究进展按正式出版的具体结果引用，不把一个方法名称当作整个猜想的证明。

\subsection{限制猜想}
\label{b2:subsec:230601}

对紧非零曲率超曲面，一种经典延拓目标是
\[
 E:L^s(S)\to L^q(\mathbb R^d),\qquad
 q>\frac{2d}{d-1},\quad
 \frac{d+1}{q}\le\frac{d-1}{s'}.
\]
必要条件已在本章证明；称为限制猜想时，这些不等式描述的是目标范围，不是本章已经取得的定理。具体维数、曲面及端点版本需要分别陈述。本书采用这一经典目标作为组织方法的背景，不以未核验的“最新范围”替代准确的定理。

\ProseParagraphBreak
当 \(s=2\)，Stein–Tomas 已达到帽必要条件边界。进一步的困难主要在输入指数改变后的范围，以及许多方向相互作用的控制。球面、抛物面与锥面有不同对称性和曲率结构，不能无说明地共用同一端点命题。

\subsection{双线性归约}
\label{b2:subsec:230602}

把 \(g\) 分成小帽后，\(|Eg|^2\) 是不同帽延拓的乘积之和。
相邻帽的乘积可在更细尺度继续分解；分离帽的法向有角度，则有额外双线性增益。这形成双线性归约的基本思路。
归约的完整形式需控制近对角项、远对角项、尺度求和及输出空间的准范数，不是简单地丢弃相邻帽。

\ProseParagraphBreak
\chapref{b2:ch:25}将完整计算二维抛物线上分离区间的乘积 \(L^2\) 估计：和映射的 Jacobi 行列式恰好等于法向分离量。这个计算是现代双线性思想的一个可独立核验模型。一般高维双线性限制理论需要更复杂的几何，不作为本章核心证明依赖。

\subsection{波包}
\label{b2:subsec:230603}

一个半径 \(R^{-1/2}\) 的曲面帽产生切向宽度 \(R^{1/2}\)、法向长度 \(R\) 的波包。先在帽内做频率局部化，再按切向空间位置分解输入，可把延拓近似组织成许多这样的管。
频率宽度与空间宽度互为倒数，法向长度则来自曲率的二次变化，这与 Knapp 管的计算一致。

\ProseParagraphBreak
波包不是严格紧支撑于一条管：光滑截断提供的是快速衰减尾部。将波包当作管特征函数会丢失相位、正交性和尾部误差。完整方法必须同时记录系数平方和、管的方向及误差求和。

\subsection{解耦}
\label{b2:subsec:230604}

解耦比较总延拓的 \(L^q\) 范数与各个小帽延拓范数的平方和。
它允许各帽内部保持相干，再控制帽与帽之间的组合；因此不同于只对输入使用 Plancherel。
抛物面的一个正式出版结果见 \cite{BourgainDemeter2015Proof}，系统处理见 \cite{Demeter2020Fourier}；\chapref{b2:ch:25}将准确写出其带尺度损失的模型版本并标注深层输入。

\ProseParagraphBreak
不同曲面、不同切分尺度和不同 \(q\) 对应不同解耦常数。必须保留加权空间球、帽大小与 \(R^\epsilon\) 因子，才能在尺度归纳中使用。把解耦粗略概括成“帽之间正交”会把 \(L^2\) 的等距事实与非 Hilbert 空间的深层估计混淆。

\subsection{多项式分割}
\label{b2:subsec:230605}

多项式分割把空间中的质量分配到许多胞腔，并把靠近多项式零集的波包单独处理。横穿胞腔的管可利用交点次数，贴近零集的管则需低维几何或横截估计。这两部分的不同控制是方法的核心。
相关正式出版的限制应用见 \cite{Guth2016Restriction}。本章只说明方法处理什么障碍，不在没有完整胞腔计数、切向管分析和尺度递推的情况下宣称一个新范围。

\ProseParagraphBreak
现代限制理论的共同组织原则是先辨认局部相干的波包，再利用方向之间的几何分离。驻相、平方函数、横截性和尺度归纳分别处理不同环节。阅读具体定理时应先核对曲面、输入指数、输出指数、局部半径和允许的尺度损失，再判断它能否用于目标问题。


\begin{chaptersummary}
限制定理的对象是一个从测试函数连续延拓的算子。曲面测度、环境测度和带权输入范数承担不同职责；改变曲面尺度或密度时，相关因子必须重新计算。延拓对偶和正二次型使这些空间之间的关系可以核对，但二次型非负并不意味着卷积核处处非负，也不意味着算子保持函数的点态非负性。

\ProseParagraphBreak

帽集中给出一条可取等号的指数边界，球面尾部则要求另一个严格不等式。Stein--Tomas 在平方曲面输入线上达到帽边界；它的证明保留了曲面结构与核的振荡，不能由取绝对值后的衰减估计替代。径向输入、加权空间或不同频率测度可以改变问题，因此使用结论前应先确定这些附加结构。

\ProseParagraphBreak

最佳常数是否达到，需要有界性之外的紧性。调制、帽集中与多个轮廓的分离是不同损失方式，处理一种并不自动排除其余方式。现代方法用波包及方向几何组织大量帽的相互作用；管形描述还须配合相位、系数和尾部估计。单个模型的计算能够说明尺度，却不能独自证明一般限制猜想。
\end{chaptersummary}

\BookExercises

% 自拟；基础到中等：平面迹的最优加权补偿，求准确常数并证明临界权重失败。
\Needspace{6\baselineskip}
\begin{exercise}\label{b2:ex:23-new-1}
设整数 \(d\geq2\)、\(s\geq0\)，把 \(x\in\mathbb R^d\) 写成 \((x',t)\)。对 \(f\in\mathcal S(\mathbb R^d)\)，令
\[
 R_0f(\xi')=\widehat f(\xi',0),
 \quad
 \|f\|_{H_s}^2=\int_{\mathbb R^d}(1+t^2)^s\cdot|f(x',t)|^2\,\mathrm dx'\,\mathrm dt.
\]
这里 \(H_s\) 表示使上述范数有限的可测函数等价类空间。证明存在有限常数 \(C_s\)，使
\[
 \|R_0f\|_{L^2(\mathbb R^{d-1})}\leq C_s\|f\|_{H_s}
\]
当且仅当 \(s>1/2\)，并在成立时求最佳常数
\[
 C_s=\left(\int_{\mathbb R}(1+t^2)^{-s}\,\mathrm dt\right)^{1/2}.
\]
说明所得延拓怎样定义于全部 \(H_s\)，以及权重如何补偿平面法向的无限相干。
\end{exercise}
\begin{solution}
先对 Schwartz 输入用 Fubini，得到 \(R_0f\) 是函数
\(F(x')=\int_{\mathbb R}f(x',t)\,\mathrm dt\) 关于 \(x'\) 的 Fourier 变换。因此由切向 Plancherel 与加权 Cauchy--Schwarz，
\[
 \begin{aligned}
 \|R_0f\|_2^2
 &=\int\left|\int f(x',t)\,\mathrm dt\right|^2\,\mathrm dx'\\
 &\leq\left(\int(1+t^2)^{-s}\,\mathrm dt\right)
       \int(1+t^2)^s\cdot|f(x',t)|^2\,\mathrm dx'\,\mathrm dt.
 \end{aligned}
\]
第一个积分有限恰好要求 \(2s>1\)。在此条件下，同一个不等式保证一般 \(H_s\) 输入的法向积分对几乎每个 \(x'\) 绝对收敛，并定义一个 \(L^2(\mathbb R^{d-1})\) 函数。对它作平方意义的 Fourier 变换，就是所需延拓；Schwartz 类在加权平方范数中稠密，故该延拓唯一。

为证明常数和必要性，取 \(u\in\mathcal S(\mathbb R^{d-1})\)、\(\|u\|_2=1\)，考虑乘积输入 \(u(x')h(t)\)。记
\[
 I_R=\int_{-R}^R(1+t^2)^{-s}\,\mathrm dt,
 \quad h_R(t)=(1+t^2)^{-s}\cdot\mathbf1_{[-R,R]}(t).
\]
则 \(\int h_R=I_R\)，而 \(\int(1+t^2)^s|h_R|^2=I_R\)，所以相应的输出与输入范数之比为 \(I_R^{1/2}\)。虽然 \(h_R\) 未必光滑，但可在一个固定的稍大紧区间内用 \(C_c^\infty\) 函数逼近；该区间上权重有正上下界，平方逼近也使积分收敛。因此 Schwartz 输入可以任意接近同一比值。

若 \(s>1/2\)，令 \(R\to\infty\)，下界趋于所给 \(C_s\)，与上界一致。若 \(s\leq1/2\)，则 \(I_R\to\infty\)，不可能有统一的估计；其中 \(s=1/2\) 时 \(I_R\) 按 \(\log R\) 增长。这表明临界失败虽缓慢，却不能由增加一个固定常数修复。权重所提供的正是法向积分泛函的有界性，整个论证并未使用曲率。
\end{solution}

% 自拟；中等：可逆仿射变换下的延拓范数，区分推前测度与椭球面积测度的非恒定Jacobian。
\Needspace{6\baselineskip}
\begin{exercise}\label{b2:ex:23-new-2}
设整数 \(d\geq2\)，\(S=S^{d-1}\) 取面积测度 \(\sigma\)，固定
\(2(d+1)/(d-1)\leq q<\infty\)。记球面延拓 \(E:L^2(S)\to L^q(\mathbb R^d)\) 的最佳常数为 \(C_q\)。设 \(A\) 为可逆实矩阵、\(b\in\mathbb R^d\)，令
\[
 F(\xi)=A\xi+b,\quad \Sigma=F(S),\quad D=|\det A|.
\]
在 \(\Sigma\) 上先取推前测度 \(\mu=F_\#\sigma\)，即
\(\int H\,\mathrm d\mu=\int_S H(F(\xi))\,\mathrm d\sigma(\xi)\)。证明相应延拓的最佳 \(L^2(\mu)\to L^q\) 常数恰为 \(D^{-1/q}C_q\)。

\ProseParagraphBreak

再改用 \(\Sigma\) 的面积测度，证明面积因子为
\[
 J_A(\xi)=D|A^{-\mathsf T}\xi|.
\]
若 \(J_{\min}\)、\(J_{\max}\) 是它在球面上的最小值与最大值，证明此时的最佳常数 \(C_{\Sigma,q}\) 满足
\[
 D^{-1/q}J_{\min}^{1/2}C_q
 \leq C_{\Sigma,q}
 \leq D^{-1/q}J_{\max}^{1/2}C_q.
\]
对 \(A=\operatorname{diag}(a_1,\ldots,a_d)\)、\(a_j>0\)，明确写出两个面积因子。这说明只记录行列式是否足以处理原生面积测度。
\end{exercise}
\begin{solution}
令 \(h=g\circ F\)。推前测度的定义给出 \(\|h\|_{L^2(S)}=\|g\|_{L^2(\mu)}\)，且
\[
 E_\mu g(x)=e^{2\pi\mathrm i x\cdot b}\cdot Eh(A^{\mathsf T}x).
\]
换元 \(y=A^{\mathsf T}x\) 后，输出范数等于 \(D^{-1/q}\|Eh\|_q\)。输入的对应是满射等距映射，因而不仅有上界，两个最佳常数也恰好相差 \(D^{-1/q}\)。

在 \(\xi\in S\) 处取一组正交切向单位向量 \(v_1,\ldots,v_{d-1}\)。变换后面积是 \(Av_1,\ldots,Av_{d-1}\) 张成的平行多面体体积。其法向面积向量按余因子矩阵变换，等于 \((\det A)A^{-\mathsf T}\xi\)，至多差一个取向符号；取长度即得 \(J_A\)。也可将任意向量补为最后一列，比较两个行列式，得到同一余因子公式。因此
\[
 \int_\Sigma H(\eta)\,\mathrm d\sigma_\Sigma(\eta)
 =\int_S H(F(\xi))\cdot J_A(\xi)\,\mathrm d\sigma(\xi).
\]

令 \(h=(g\circ F)\sqrt{J_A}\)，则 \(\|h\|_2=\|g\|_{L^2(\Sigma)}\)，而面积延拓变为
\(e^{2\pi\mathrm i x\cdot b}\cdot E(\sqrt{J_A}\,h)(A^{\mathsf T}x)\)。记乘法算子 \(M h=\sqrt{J_A}\,h\)。于是
\[
 C_{\Sigma,q}=D^{-1/q}\|EM\|_{2\to q},
 \quad \|M\|_{2\to2}=J_{\max}^{1/2},
 \quad \|M^{-1}\|_{2\to2}=J_{\min}^{-1/2}.
\]
由 \(EM\) 的复合上界及 \(E=(EM)M^{-1}\)，分别得到题中上界和下界。

对正对角矩阵，\(|A^{-\mathsf T}\xi|\) 在单位球上的最小值为 \(1/\max_j a_j\)，最大值为 \(1/\min_j a_j\)。因此
\[
 J_{\min}=\frac D{\max_j a_j},
 \quad J_{\max}=\frac D{\min_j a_j}.
\]
各向异性变形使面积因子随法向改变，行列式本身没有记录这项变化。若所有 \(a_j=R\)，两个因子才都化为 \(R^{d-1}\)，并恢复球面缩放的准确幂次。
\end{solution}

% 自拟；综合应用：薄球壳频率投影的两种算子范数及其准确平方关系，连接限制、余面积与伴随。
\Needspace{6\baselineskip}
\begin{exercise}\label{b2:ex:23-new-3}
设整数 \(d\geq2\)、\(1\leq p\leq p_*=2(d+1)/(d+3)\)，对 \(0<\varepsilon<1/4\) 令
\[
 \Omega_\varepsilon=\{\,\xi\in\mathbb R^d:1-\varepsilon<|\xi|<1+\varepsilon\,\},
 \quad
 P_\varepsilon f=\mathcal F^{-1}(\mathbf1_{\Omega_\varepsilon}\widehat f).
\]
算子先在 Schwartz 类上定义。证明它有相应的连续延拓，并有
\[
 \|P_\varepsilon\|_{L^p\to L^2}\asymp\varepsilon^{1/2},
 \quad
 \|P_\varepsilon\|_{L^p\to L^{p'}}
   =\|P_\varepsilon\|_{L^p\to L^2}^2
   \asymp\varepsilon.
\]
隐含常数只依赖 \(d,p\)。这里不要求证明 \(P_\varepsilon\) 在同一个 \(L^p\) 空间上有界。
\end{exercise}
\begin{solution}
Plancherel、极坐标和半径缩放后的 Stein--Tomas 估计给出
\[
 \begin{aligned}
 \|P_\varepsilon f\|_2^2
 &=\int_{1-\varepsilon}^{1+\varepsilon}
       \|\widehat f\|_{L^2(rS^{d-1})}^2\,\mathrm dr\\
 &\leq C_{d,p}\|f\|_p^2
       \int_{1-\varepsilon}^{1+\varepsilon}
              r^{d-1-2d/p'}\,\mathrm dr
 \leq C'_{d,p}\varepsilon\|f\|_p^2.
 \end{aligned}
\]
半径始终位于 \((3/4,5/4)\)，所以幂因子一致有界。有限 \(p\) 下的稠密性给出 \(L^p\to L^2\) 延拓。为求下界，固定 \(\widehat f\in C_c^\infty\)，使其在 \(1/2\leq|\xi|\leq3/2\) 上恒为一。这个 \(f\) 是与 \(\varepsilon\) 无关的 Schwartz 函数，而
\[
 \|P_\varepsilon f\|_2^2=|\Omega_\varepsilon|\asymp\varepsilon.
\]
除以固定的 \(\|f\|_p^2\)，即得第一条范数等价。

把已证的 \(L^p\to L^2\) 算子记为 \(B\)。其伴随 \(B^*:L^2\to L^{p'}\) 有同一范数，包括 \(p=1\) 时的 \(L^\infty\) 值域。在测试函数上，实值幂等符号给出
\[
 \langle Bf,Bh\rangle_{L^2}
 =\int_{\Omega_\varepsilon}\widehat f(\xi)\cdot\overline{\widehat h(\xi)}\,\mathrm d\xi
 =\langle P_\varepsilon f,h\rangle.
\]
因此 \(B^*B\) 与原来 \(P_\varepsilon\) 的分布定义一致，提供 \(L^p\to L^{p'}\) 延拓，并满足 \(\|P_\varepsilon\|_{p\to p'}\leq\|B\|^2\)。反过来，对测试输入有
\[
 \|Bf\|_2^2
 =\langle P_\varepsilon f,f\rangle
 \leq\|P_\varepsilon\|_{p\to p'}\|f\|_p^2.
\]
对单位测试输入取上确界，得到反向不等式，故两范数恰有平方关系。这个论证利用了中间空间的 Hilbert 结构；第二条估计并非把第一条的输出指数直接换掉。
\end{solution}

% 自拟；基础核验与反例：显式球面核区分非负二次型与保持点态非负，并识别二次型的非平凡零空间。
\Needspace{6\baselineskip}
\begin{exercise}\label{b2:ex:23-new-4}
在 \(\mathbb R^3\) 中令 \(\sigma\) 为单位球面的面积测度，\(K=\check\sigma\)，并对 \(f\in\mathcal S(\mathbb R^3)\) 置 \(Tf=K*f\)。不用 Bessel 函数，直接证明
\[
 K(x)=\frac{2\sin(2\pi|x|)}{|x|}\quad(x\ne0),
 \quad K(0)=4\pi.
\]
证明 \(\int Tf\cdot\overline f\geq0\)，同时构造非负 \(f\in C_c^\infty\)，使 \(Tf\) 在某个开集内为负。再构造非零 Schwartz 函数满足 \(Tf=0\)。解释这三项性质为何可以同时成立。
\end{exercise}
\begin{solution}
由旋转对称性，可令 \(x=(0,0,r)\)。用球面极坐标并作换元 \(u=\cos\theta\)，得到
\[
 K(x)=2\pi\int_0^\pi e^{2\pi\mathrm i r\cos\theta}\sin\theta\,\mathrm d\theta
 =2\pi\int_{-1}^1e^{2\pi\mathrm i ru}\,\mathrm du
 =\frac{2\sin(2\pi r)}r.
\]
原点极限为 \(4\pi\)，也等于球面面积。对 Schwartz 输入，\(f\in L^1\) 且 \(K\) 有界，下面的积分换序绝对可行：
\[
 \int_{\mathbb R^3}Tf(x)\cdot\overline{f(x)}\,\mathrm dx
 =\int_{S^2}|\widehat f(\xi)|^2\,\mathrm d\sigma(\xi)\geq0.
\]
这里非负的是积分二次型。

取 \(x_0=(0,0,3/4)\)，则 \(K(x_0)<0\)。由连续性，存在 \(\delta>0\)，使 \(K\) 在 \(B(x_0,2\delta)\) 上为负。选非负 \(f\in C_c^\infty(B(0,\delta))\)，满足 \(\int f=1\)。对 \(x\in B(x_0,\delta)\) 和 \(y\in\operatorname{supp}f\)，有 \(x-y\in B(x_0,2\delta)\)，故
\(Tf(x)=\int K(x-y)\cdot f(y)\,\mathrm dy<0\)。所以这个算子不保持函数的点态非负性。

最后取非零 \(b\in C_c^\infty(B(0,1/2))\)，令 \(f=\check b\)。此时 \(f\in\mathcal S\setminus\{0\}\)，但 \(\widehat f|_{S^2}=0\)，并且
\[
 Tf(x)=\int_{S^2}e^{2\pi\mathrm i x\cdot\xi}\widehat f(\xi)\,\mathrm d\sigma(\xi)=0.
\]
二次型只测量 Fourier 变换在球面上的部分，既不检测所有非零输入，也不要求空间核取非负值。因而它的非负性、非平凡零空间和点态变号没有冲突。
\end{solution}

% 自拟；综合证明：径向输入的准确限制范围，由径向对偶求最佳常数并量化临界对数发散。
\Needspace{6\baselineskip}
\begin{exercise}\label{b2:ex:23-new-5}
设整数 \(d\geq2\)、\(1\leq p<\infty\)，只对径向 Schwartz 函数 \(f\) 考虑球面限制。记
\[
 A_d=\sigma(S^{d-1}),
 \quad j(x)=A_d^{-1}\int_{S^{d-1}}e^{2\pi\mathrm i x\cdot\xi}\,\mathrm d\sigma(\xi).
\]
证明存在统一的限制估计
\[
 \|\widehat f|_{S^{d-1}}\|_2\leq C_{d,p}\|f\|_p
\]
当且仅当 \(1\leq p<2d/(d+1)\)。成立时，证明径向类中的最佳常数是 \(A_d^{1/2}\|j\|_{p'}\)。在临界指数 \(p=2d/(d+1)\)，构造径向单位输入，其限制范数趋于无穷；可用截断对偶函数获得至少 \((\log R)^{(d-1)/(2d)}\) 的增长。
\end{exercise}
\begin{solution}
径向性使 \(\widehat f\) 在单位球面上为常数。对球面方向平均，并使用 Fubini，得到
\[
 \widehat f(e_1)=A_d^{-1}\int_{S^{d-1}}\widehat f(\xi)\,\mathrm d\sigma(\xi)
 =\int_{\mathbb R^d}f(x)\cdot j(x)\,\mathrm dx.
\]
球面对称使 \(j\) 为实值径向函数，因此最后一式没有反射或共轭差异。限制范数便是 \(A_d^{1/2}\) 乘这个积分的绝对值，Hölder 给出所述常数作为上界。

对 \(1<p<\infty\)，令 \(q=p'\)。在任意球 \(B(0,R)\) 上，径向函数
\(\operatorname{sgn}(j)|j|^{q-1}\cdot\mathbf1_{B(0,R)}\) 是这一积分泛函的对偶取等函数；零点处取零。归一化后，输入 \(L^p\) 范数为一，积分值为 \(\|j\|_{L^q(B(0,R))}\)。这些紧支撑径向函数可用径向 \(C_c^\infty\) 函数在 \(L^p\) 中逼近，例如用径向光滑核卷积。逼近支撑保持在一个稍大球内，那里 \(j\in L^q\)，所以积分值也收敛。再归一化并令 \(R\to\infty\)，就证明了最佳常数公式；若 \(j\notin L^q\)，同一构造也证明无界。

正文的球面驻相展开给出
\[
 |j(x)|\leq C_d(1+|x|)^{-(d-1)/2},
\]
并在每个充分大的单位径向区间中，找到占固定比例的一部分，使 \(|j(r)|\geq c_dr^{-(d-1)/2}\)。因此对有限 \(q\)，
\[
 j\in L^q(\mathbb R^d)
 \quad\Longleftrightarrow\quad q>\frac{2d}{d-1}.
\]
转回 \(p\) 就是 \(p<2d/(d+1)\)。当 \(p=1\) 时，\(|j|\leq1\) 且 \(j(0)=1\)，故 \(\|j\|_\infty=1\)；取向原点集中的非负径向光滑近似恒等族，积分趋于一，仍得到相同的最佳常数公式。

在临界 \(q_c=2d/(d-1)\) 处，上述上下界还给出
\[
 \int_{|x|<R}|j(x)|^{q_c}\,\mathrm dx\asymp\log R
 \quad(R\geq2).
\]
用前面的截断取等函数，选择其径向光滑逼近，使归一化后的积分至少保留一半，便得到所求径向 Schwartz 单位输入，限制范数至少为 \(c_d(\log R)^{1/q_c}\)。径向假设禁止任意单帽集中，所以必要条件由振荡尾部决定；这解释了本题范围为何能够超过一般输入的 Stein--Tomas 范围。
\end{solution}

% 自拟；中等：紧频率支集的整函数延拓与开集唯一性，给出不依赖PDE唯一延拓定理的证明。
\Needspace{6\baselineskip}
\begin{exercise}\label{b2:ex:23-new-6}
设整数 \(d\geq2\)，\(S\subset\mathbb R^d\) 为非空紧致光滑超曲面，取面积测度，并设 \(g\in L^1(S)\)。令 \(M=\max_{\xi\in S}|\xi|\)。证明延拓积分在 \(z\in\mathbb C^d\) 上仍有定义，并满足
\[
 |Eg(z)|\leq e^{2\pi M|\operatorname{Im}z|}\|g\|_1,
 \quad
 \|\partial^\alpha Eg\|_{L^\infty(\mathbb R^d)}
       \leq(2\pi M)^{|\alpha|}\|g\|_1.
\]
证明它在每一点有全局收敛的复幂级数。由此推出：若 \(Eg\) 在一个非空实开集上恒零，则 \(g=0\) 几乎处处；特别地，非零数据不能产生具有紧空间支集的延拓函数。不得直接调用偏微分方程的唯一延拓定理。
\end{exercise}
\begin{solution}
由于
\(\left|e^{2\pi\mathrm i z\cdot\xi}\right|
=e^{-2\pi\operatorname{Im}z\cdot\xi}\leq e^{2\pi M|\operatorname{Im}z|}\)，
积分绝对收敛，并立刻给出第一条估计。固定 \(z_0\in\mathbb C^d\)，把指数写成
\[
 e^{2\pi\mathrm i z\cdot\xi}
 =e^{2\pi\mathrm i z_0\cdot\xi}
   \sum_{k=0}^\infty\frac{\bigl(2\pi\mathrm i(z-z_0)\cdot\xi\bigr)^k}{k!}.
\]
当 \(|z-z_0|\leq R\) 时，级数绝对值之和至多为
\(e^{2\pi M|\operatorname{Im}z_0|}e^{2\pi MR}\)，乘 \(|g|\) 后可积。因此可以逐项积分，在每个复紧球上一致收敛；再把各次幂按多重指标展开，便得到在整个 \(\mathbb C^d\) 收敛的幂级数。逐项求导或在原积分中求导，得到
\[
 \partial^\alpha Eg(x)
 =\int_S(2\pi\mathrm i\xi)^\alpha
             e^{2\pi\mathrm i x\cdot\xi}g(\xi)\,\mathrm d\sigma(\xi).
\]
对实 \(x\) 取绝对值，就有第二条估计。

若 \(Eg\) 在非空实开集 \(U\) 上为零，取 \(x_0\in U\)。它在 \(x_0\) 的所有实偏导数都为零，因而上面以 \(z_0=x_0\) 展开的全部幂级数系数都为零。全局收敛说明 \(Eg\) 在整个实空间中恒零。

最后须由输出为零回到输入。对任意 \(\varphi\in\mathcal S(\mathbb R^d)\)，Fourier 反演与 Fubini 给出
\[
 \int_{\mathbb R^d}Eg(x)\cdot\widehat\varphi(x)\,\mathrm dx
 =\int_S\varphi(\xi)g(\xi)\,\mathrm d\sigma(\xi).
\]
绝对换序由 \(\|g\|_1\cdot\|\widehat\varphi\|_1<\infty\) 保证。左边为零，说明有限测度 \(g\,\mathrm d\sigma\) 作为分布为零，所以 \(g=0\) 几乎处处。具有紧支集的连续输出在某个非空外部开集为零，故也只能来自零数据。
\end{solution}

% 自拟；选做综合：L²到L∞延拓的全部极值与精确稳定性，补充有限输出指数下的紧性讨论。
\Needspace{6\baselineskip}
\begin{optionalexercise}\label{b2:ex:23-new-7}
设整数 \(d\geq2\)，\(S=S^{d-1}\) 取面积测度，\(A_d=\sigma(S)\)，\(E\) 为球面延拓。证明 \(E:L^2(S)\to L^\infty\) 的最佳常数是 \(A_d^{1/2}\)，且全部单位范数极值函数恰为
\[
 g_{a,\theta}(\xi)=A_d^{-1/2}e^{\mathrm i\theta}e^{-2\pi\mathrm i a\cdot\xi},
 \quad a\in\mathbb R^d,\quad\theta\in\mathbb R.
\]
进一步证明，对每个 \(\|g\|_2=1\)，
\[
 \inf_{a,\theta}\|g-g_{a,\theta}\|_2^2
 =2\left(1-\frac{\|Eg\|_\infty}{A_d^{1/2}}\right),
\]
且此下确界能取到。特别地，若 \(\|Eg\|_\infty\geq(1-\varepsilon)A_d^{1/2}\)，则 \(g\) 距离某个上述极值函数至多 \(\sqrt{2\varepsilon}\)。需要说明输出为何确实在某一点达到其最大模。
\end{optionalexercise}
\begin{solution}
Cauchy--Schwarz 给出 \(|Eg(x)|\leq A_d^{1/2}\|g\|_2\)。所列 \(g_{a,\theta}\) 的平方范数为一，并满足
\(Eg_{a,\theta}(a)=A_d^{1/2}e^{\mathrm i\theta}\)，所以最佳常数恰为 \(A_d^{1/2}\)。

对一般 \(g\in L^2(S)\)，先用光滑函数 \(g_n\) 在 \(L^2(S)\) 中逼近。光滑球面密度的驻相估计说明 \(Eg_n(x)\to0\) 当 \(|x|\to\infty\)，而
\[
 \|Eg_n-Eg\|_\infty\leq A_d^{1/2}\|g_n-g\|_2\longrightarrow0.
\]
因此 \(Eg\) 也是连续且在无穷远趋零的函数。若 \(g\ne0\)，Fourier 变换的单射性保证 \(Eg\) 不恒零；选一个模为正的值，充分大的球外模小于其一半，故最大模在一个紧球内达到。

若单位输入是极值函数，在某个最大点 \(a\) 有 \(|Eg(a)|=A_d^{1/2}\)。Cauchy--Schwarz 的取等条件迫使 \(g\) 与 \(e^{-2\pi\mathrm i a\cdot\xi}\) 成比例，单位范数再确定比例的模，恰得到 \(g_{a,\theta}\)。这证明了分类。

为求距离，令 \(h_a=A_d^{-1/2}e^{-2\pi\mathrm i a\cdot\xi}\)。复内积在第一变量线性，故
\[
 \langle g,h_a\rangle=A_d^{-1/2}Eg(a),
 \quad
 \|g-e^{\mathrm i\theta}h_a\|_2^2
 =2-2\operatorname{Re}\bigl(e^{-\mathrm i\theta}A_d^{-1/2}Eg(a)\bigr).
\]
固定 \(a\)，选相位使实部等于模，得到最小值
\(2-2A_d^{-1/2}|Eg(a)|\)。再选达到输出最大模的 \(a\)，就得到准确公式，并说明下确界取到。最后代入假设并开平方，即得稳定性界。这个极值问题由单点 Cauchy--Schwarz 取等完全决定，与有限输出指数时必须控制整个空间质量的情形有所不同。
\end{solution}

% 自拟；选做综合：Gaussian抛物面波包的全时空显式计算，核对管状尺度、长时间扩散与可积阈值。
\Needspace{6\baselineskip}
\begin{optionalexercise}\label{b2:ex:23-new-8}
设整数 \(n\geq1\)。在全抛物面 \(\{\,(\eta,|\eta|^2):\eta\in\mathbb R^n\,\}\) 上使用参数测度 \(\mathrm d\eta\)，定义
\[
 \mathcal Eg(x,t)=\int_{\mathbb R^n}
          e^{2\pi\mathrm i(x\cdot\eta+t|\eta|^2)}g(\eta)\,\mathrm d\eta.
\]
对 \(\delta>0\)、\(\eta_0\in\mathbb R^n\)，取单位平方范数数据
\[
 g_{\delta,\eta_0}(\eta)=2^{n/4}\delta^{-n/2}
              \cdot e^{-\pi|\eta-\eta_0|^2/\delta^2}.
\]
求其延拓的显式公式，并证明模的空间宽度随时间为
\[
 W_\delta(t)=\delta^{-1}\sqrt{1+4t^2\delta^4},
\]
中心位于 \(x=-2t\eta_0\)。对 \(1\leq q<\infty\)，证明
\[
 \mathcal Eg_{\delta,\eta_0}\in L^q(\mathbb R^{n+1})
 \quad\Longleftrightarrow\quad q>\frac{2(n+1)}n,
\]
且在成立时
\[
 \|\mathcal Eg_{\delta,\eta_0}\|_q
 =C_{n,q}\delta^{n/2-(n+2)/q},\quad0<C_{n,q}<\infty.
\]
解释 \(|t|\lesssim\delta^{-2}\) 时的管状描述，以及更长时间为何不能保持同一管宽。这里的模型不是带面积测度的紧球面。
\end{optionalexercise}
\begin{solution}
换元 \(\eta=\eta_0+\delta u\)，首先验证
\[
 \|g_{\delta,\eta_0}\|_2^2
 =2^{n/2}\int_{\mathbb R^n}e^{-2\pi|u|^2}\,\mathrm du=1.
\]
同一换元将延拓化为一个复 Gaussian 积分。对于 \(\operatorname{Re}a>0\)，有
\[
 \int_{\mathbb R^n}e^{-\pi a|u|^2}e^{2\pi\mathrm i v\cdot u}\,\mathrm du
 =a^{-n/2}e^{-\pi|v|^2/a}\quad(v\in\mathbb R^n).
\]
这里 \(a^{-n/2}\) 取右半平面上在正实轴为正的解析分支。该式先由实 Gaussian 的 Fourier 公式对 \(a>0\) 得到；参数位于右半平面紧集时，积分及参数导数均被一个多项式乘固定 Gaussian 控制，故两侧全纯，恒等定理把等式推广到整个右半平面。

取 \(a=1-2\mathrm i t\delta^2\)、\(v=\delta(x+2t\eta_0)\)，得到
\[
 \begin{aligned}
 \mathcal Eg_{\delta,\eta_0}(x,t)
 &=2^{n/4}\delta^{n/2}
   \cdot e^{2\pi\mathrm i(x\cdot\eta_0+t|\eta_0|^2)}\\
 &\quad\cdot(1-2\mathrm i t\delta^2)^{-n/2}
   \cdot\exp\left(-\frac{\pi\delta^2|x+2t\eta_0|^2}
                       {1-2\mathrm i t\delta^2}\right).
 \end{aligned}
\]
记 \(B(t)=1+4t^2\delta^4\)。取绝对值得
\[
 |\mathcal Eg_{\delta,\eta_0}(x,t)|
 =2^{n/4}\delta^{n/2}B(t)^{-n/4}
       \exp\left(-\frac{\pi\delta^2|x+2t\eta_0|^2}{B(t)}\right).
\]
指数中的二次项给出题述中心和宽度。当 \(|t|\leq c\delta^{-2}\) 时，\(B(t)\) 有依赖 \(c\) 的统一上下界，宽度与 \(\delta^{-1}\) 可比，中心模与 \(\delta^{n/2}\) 可比；这就是有限时间内沿倾斜中心线的一根管。若 \(|t|\gg\delta^{-2}\)，宽度则与 \(\delta|t|\) 可比，不能继续用固定宽度描述整个包。Gaussian 在任何空间点都不严格为零，管外描述的是衰减而非支集。

最后直接积分，避免把单个时刻的界误当成全时空界。对每个 \(t\)，平移中心并作 Gaussian 换元，得到
\[
 \int_{\mathbb R^n}|\mathcal Eg_{\delta,\eta_0}(x,t)|^q\,\mathrm dx
 =2^{nq/4}q^{-n/2}\delta^{nq/2-n}
       (1+4t^2\delta^4)^{-n(q-2)/4}.
\]
再令 \(u=t\delta^2\)，Tonelli 给出
\[
 \|\mathcal Eg_{\delta,\eta_0}\|_q^q
 =2^{nq/4}q^{-n/2}\delta^{nq/2-n-2}
       \int_{\mathbb R}(1+4u^2)^{-n(q-2)/4}\,\mathrm du.
\]
末项有限当且仅当 \(n(q-2)/2>1\)，即 \(q>2(n+1)/n\)；等号时仍有对数发散。积分有限时开 \(q\) 次方，就是所给准确尺度。特别在 \(q=2(n+2)/n\) 时，\(\delta\) 的幂为零，与平方输入的抛物面尺度相符；这一单个 Gaussian 家族的计算本身没有证明对全部输入的延拓估计。
\end{solution}
